\providecommand{\oritoneexamplemode}{white}
\documentclass[\oritoneexamplemode,11pt]{oritone-manuscript}

\title{A Complete Manuscript Example for Reusable Scientific Reports}
\author[1]{Ada Example}
\author[1,2]{Benoit Sample}
\author[2]{Clara Rivera}
\affil[1]{Department of Applied Typesetting, Example University}
\affil[2]{Institute for Reusable Documents, Sample Laboratory}
\date{\today}

\printcompileinfo

\begin{document}
\maketitle

\begin{abstract}
This document demonstrates the manuscript class on content that resembles a
short technical paper rather than a minimal smoke test. It includes authors and
affiliations, keywords, line numbers, compact headings, mathematical notation,
figures, tables, listings, cross-references, links, captions, and enough body
text to show running heads and filled pages.
\end{abstract}

\keywords{document classes\keywordsep scientific manuscripts\keywordsep reusable LaTeX\keywordsep
research software}

\modulolinenumbers[5]
\linenumbers

\section{Introduction}
\label{sec:introduction}

Reusable manuscript infrastructure matters when research groups write many
short reports with similar typographic expectations. A compact class should
therefore provide the routine decisions once: page geometry, line spacing,
author formatting, captions, references, links, listings, and running heads. It
should also stay quiet enough that authors can focus on structure and argument
instead of presentation.

The example is intentionally filled with ordinary paragraphs. The first page
shows the title block, abstract, keyword line, and the beginning of the paper.
Subsequent pages show how headings sit next to dense prose, floats, numbered
equations, and code. Cross-references such as \Cref{sec:methods} and
\Cref{fig:workflow} use the class defaults, while links such as
\href{https://example.com/protocol}{the project protocol} use the configured
semantic link colors.

\subsection{Design requirements}

The class is aimed at review drafts, internal reports, preprints, and small
technical notes. Those documents often need generous line spacing, visible page
state, line numbers, and compact headings. They also need predictable behavior
for standard floats. The typography should be restrained: manuscript pages
usually benefit from clear hierarchy, not elaborate decoration.

\begin{oritonenote}[Design Note]
The note environment uses the selected accent role. It is useful for review
comments, assumptions, or short implementation context that belongs with the
body text but should remain visually scannable.
\end{oritonenote}

\subsubsection{Scope}

The class does not own the scientific content. It owns the repeated
presentation layer. Authors remain free to choose theorem packages,
bibliography tools, discipline-specific notation, and journal-specific final
formatting when a manuscript leaves the internal review stage.

\paragraph{Implementation note}
The run-in paragraph heading demonstrates the fourth heading level and keeps a
short note attached to the first sentence. The following sections use longer
blocks of prose so that page headers, footers, captions, and line numbering are
visible in the compiled output.

\section{Methods}
\label{sec:methods}

Consider a scalar state \(u:\Omega\to\mathbb{R}\) and an observation operator
\(G\). A compact loss used for calibration can be written as
\begin{equation}
  \mathcal{J}(u,\theta)
  =
  \frac{1}{2}\lVert G(u)-y\rVert_{\Gamma^{-1}}^2
  + \lambda \int_{\Omega}
      \frac{1}{2}\lvert\nabla u\rvert^2 + V(u,\theta)\,\mathrm{d}x .
  \label{eq:objective}
\end{equation}
The notation in \Cref{eq:objective} is deliberately conventional. The point of
the example is to show equation spacing and numbering in context, not to define
a new model. In real reports the same space is commonly used for short
derivations, objective functions, balance laws, or estimator definitions.

\begin{figure}[tbp]
  \centering
  \begin{tikzpicture}[
    node distance=2.45cm,
    every node/.style={
      draw=darkblue,
      fill=verylightblue,
      rounded corners=2pt,
      align=center,
      minimum width=2.35cm,
      minimum height=0.9cm
    },
    arrow/.style={->, thick, draw=darkgreen}
  ]
    \node (model) {Model};
    \node[right of=model] (solve) {Solve};
    \node[right of=solve] (check) {Check};
    \node[below of=solve, yshift=9mm] (archive) {Archive};
    \draw[arrow] (model) -- (solve);
    \draw[arrow] (solve) -- (check);
    \draw[arrow] (check) |- (archive);
    \draw[arrow] (archive) -| (model);
  \end{tikzpicture}
  \caption{A workflow figure using semantic colors from the shared palette.}
  \label{fig:workflow}
\end{figure}

The workflow in \Cref{fig:workflow} is deliberately small, but it exercises the
same caption and color behavior as a larger diagram. In practice the model step
records assumptions and parameter bounds, the solve step executes the numerical
kernel, the check step compares diagnostics against tolerances, and the archive
step stores artifacts with enough metadata for later review.

\begin{oritoneexample}[Accepted Pattern]
Use example callouts for validated conventions, reproducible protocol snippets,
or short observations that should read as positive guidance rather than a
warning.
\end{oritoneexample}

\subsection{Quality measures}

The numerical summary in \Cref{tab:summary} uses a plain table so that the class
does not depend on a specialized table package. The columns report the size of a
small refinement study and a pair of quality measures. Dense manuscript tables
should remain readable without heavy rules or background color.

\begin{table}[tbp]
  \centering
  \caption{A compact numerical summary used throughout the example.}
  \label{tab:summary}
  \begin{tabular}{lrrr}
    \hline
    Case & Iterations & Relative error & Wall time (s) \\
    \hline
    Coarse & 18 & \(2.4\times10^{-2}\) & 0.8 \\
    Medium & 27 & \(6.1\times10^{-3}\) & 2.7 \\
    Fine & 41 & \(1.5\times10^{-3}\) & 9.4 \\
    Validation & 39 & \(1.8\times10^{-3}\) & 8.9 \\
    \hline
  \end{tabular}
\end{table}

\section{Implementation}
\label{sec:implementation}

Listings use the same framed component as the other Oritone document classes:
line numbers, semantic syntax color, and compact monospaced text.
\Cref{lst:quality} shows a small Python-like diagnostic routine. The code is
intentionally short, but it still demonstrates keywords, strings, comments,
numbers, and continuation lines.

\begin{lstlisting}[language=Python, caption={A short quality-control routine.}, label={lst:quality}]
def residual(values, tolerance=1.0e-8):
    """Return a normalized second-difference residual."""
    total = 0.0
    count = 0
    for left, center, right in zip(values[:-2], values[1:-1], values[2:]):
        total += abs(left - 2.0 * center + right)
        count += 1
    return total / max(count, 1) < tolerance
\end{lstlisting}

The surrounding text should look like a normal manuscript rather than a boxed
demo. A reader can scan from \Cref{eq:objective} to \Cref{fig:workflow},
\Cref{tab:summary}, and \Cref{lst:quality} without encountering a change in
visual language. That consistency is the primary purpose of the class.

\section{Results}
\label{sec:results}

The method was evaluated on four manufactured cases. The first two cases were
small enough to inspect by hand, while the remaining cases were used to check
the stability of the workflow and the visual density of full pages. The
relative errors in \Cref{tab:summary} decrease at the expected rate, and the
wall time grows smoothly with refinement.

Longer paragraphs are included here to exercise page filling. A typical report
contains limitations, exceptions, and small interpretation details that do not
belong in the figure captions. The class should make those paragraphs easy to
read: no decorative sidebars, no oversized headings, and no excessive color in
body text. The Oritone palette remains visible in links, figures, and listings
without competing with the manuscript itself.

\subsection{Robustness checks}

We repeated the workflow with perturbed initial conditions and with a coarser
archive cadence. The diagnostic routine continued to reject incomplete runs,
and the stored metadata was sufficient to reproduce all accepted cases. The
main operational issue was not numerical stability but naming discipline:
reports were easiest to review when parameter files, generated figures, and
tables used the same case identifiers.

\begin{oritonewarning}[Review Caveat]
Warning callouts use the danger role and should stay rare. They are intended for
failed checks, unresolved assumptions, or required reader action.
\end{oritonewarning}

\section{Discussion}
\label{sec:discussion}

The example demonstrates the class features in one compileable source file.
Authors can remove line numbers for final preprints, omit compile information
from the footer, add bibliography packages, or replace the toy model with a real
study. The visual defaults are intentionally conservative so that such changes
do not require rewriting the document.

The most useful test for a document class is a filled page. Sparse examples can
hide bad heading spacing, fragile floats, overbearing colors, and awkward footer
placement. This source therefore includes enough ordinary text to make those
interactions visible during development.

\section{Conclusion}

The manuscript class provides a compact, review-friendly surface for scientific
reports. It demonstrates predictable defaults for typography, captions, links,
line numbers, listings, equations, tables, and figures while leaving
discipline-specific packages to the document author.

\end{document}
